% docs/texte/05_theory/two_element.tex -- theory chapter "The Two-Element Interferometer" (Interferosim)
% Version: 0.1.1 (2026-10-12) -- figure source note, reference to appendix
% Notation as defined in texte/05_theory/notation.tex.
% ================================================================

\index{interferometer!two-element}\index{geometric delay}
\section{Geometric delay and fringes}

Figure~\ref{fig:two-element} shows two antennas separated by the
baseline $\vec b$. A plane wave from the direction $\hat s$ reaches
antenna~1 later than antenna~2 by the \emph{geometric delay}
\begin{equation}
  \tau_g = \frac{\vec b\cdot\hat s}{c} = \frac{b\cos\theta}{c} .
  \label{eq:taug}
\end{equation}
For a monochromatic signal the voltages are
$V_1 = V\cos[\omega(t-\tau_g)]$ and $V_2 = V\cos(\omega t)$. A
multiplying correlator forms $V_1V_2$ and averages over a time long
compared with $1/\omega$:
\begin{equation}
  R = \langle V_1 V_2\rangle = \frac{V^2}{2}\cos(\omega\tau_g) .
  \label{eq:fringe}
\end{equation}
As the source direction changes, $R$ oscillates: these are the
\emph{fringes}. Their angular spacing near $\theta = 90^\circ$ is
$\lambda/b$; this is the angular resolution scale of the baseline.
Appendix~\ref{app:fringe} explains the term \enquote{fringe} and the
related terms in detail.

\begin{figure}[htbp]
  \centering
  % docs/texte/05_theory/figures/two_element_interferometer.tex -- two-element interferometer (TikZ)
% Version: 0.1.1 (2026-10-12) -- signal labels below the lines (no overlap with 12 pt text)
% Redrawn in TikZ after docs/images/antennengeometrie2.png (two-element
% interferometer with multiplying correlator). Geometry: baseline b along
% the x axis, source direction s at angle theta to the baseline; the
% wavefront through antenna 2 meets the line from antenna 1 along s at P,
% so |0P| = b cos(theta) = c tau_g.
\begin{tikzpicture}[>=Stealth, font=\footnotesize]
  % ---- parameters ---------------------------------------------------
  \def\b{6.0}        % baseline length (drawing units)
  \def\th{55}        % angle theta between baseline and source direction
  \pgfmathsetmacro{\px}{\b*cos(\th)*cos(\th)}
  \pgfmathsetmacro{\py}{\b*cos(\th)*sin(\th)}
  \coordinate (A1) at (0,0);
  \coordinate (A2) at (\b,0);
  \coordinate (P)  at (\px,\py);

  % ---- source direction and wavefront ------------------------------
  \draw[thin] (A1) -- (P);
  \draw[thin,dotted] (P) -- ($(A1)!1.75!(P)$);
  \draw[->,thick] ($(A1)!1.75!(P)$) -- ($(A1)!1.95!(P)$)
        node[above right=-1pt] {$\hat{s}$};
  \draw[thin,dotted] (A2) -- ++(\th:3.6);
  \draw[thin] (A2) -- (P);                                   % wavefront
  \draw[thin] (A1) -- node[below=3pt] {$\vec{b}$} (A2);       % baseline
  % path difference along s
  \path (A1) -- (P) node[midway,sloped,above=2pt]
        {$c\,\tau_g = \vec{b}\cdot\hat{s}$};
  % angle theta
  \draw (1.5,0) arc[start angle=0,end angle=\th,radius=1.5];
  \node at ({\th/2}:1.8) {$\theta$};
  % right angle at P
  \draw[thin] ($(P)!0.22cm!(A1)$) -- ++(\th-90:0.22cm) -- ($(P)!0.22cm!(A2)$);

  % ---- antennas: reflector arc on a mount (triangle) ---------------
  \foreach \A/\n in {A1/1, A2/2}{
    \draw[very thick] ($(\A)+(-1.1,0.55)$) .. controls ($(\A)+(-0.45,-0.05)$)
          and ($(\A)+(0.45,-0.05)$) .. ($(\A)+(1.1,0.55)$);
    \draw[thick] (\A) -- ++(-0.38,-1.25) -- ++(0.76,0) -- cycle;
    \node at ($(\A)+(0,-0.85)$) {\n};
  }

  % ---- signal path: multiplier and averaging -----------------------
  \coordinate (M) at ({\b/2},-2.4);
  \draw[thick] (A1) ++(0,-1.25) |- ($(M)+(-0.32,0)$);
  \draw[thick] (A2) ++(0,-1.25) |- ($(M)+(0.32,0)$);
  \draw[thick] (M) circle[radius=0.32];
  \draw[thick] (M) ++(-0.2,-0.2) -- ++(0.4,0.4) (M) ++(-0.2,0.2) -- ++(0.4,-0.4);
  \node[anchor=north east] at ($(M)+(-0.4,-0.08)$) {$V_1 = V\cos[\omega(t-\tau_g)]$};
  \node[anchor=north west] at ($(M)+(0.4,-0.08)$)  {$V_2 = V\cos(\omega t)$};
  \node[draw,thick,minimum width=1.1cm,minimum height=0.55cm] (avg)
        at ($(M)+(0,-2.1)$) {$\langle\;\rangle$};
  \draw[thick] ($(M)+(0,-0.32)$) -- node[right,pos=0.6] {$V_1V_2$} (avg.north);

  % ---- correlator output: fringe pattern ----------------------------
  \draw[->,thick] (avg.east) -- ++(0.5,0);
  \draw[thin,domain=0:6,samples=400,smooth]
     plot ({\b/2+1.25+\x},{-4.5+0.42*exp(-((\x-3)^2)/1.8)*cos(deg(9.5*\x))});
  \node at ({\b/2+4.25},-5.2) {$R = \tfrac{V^2}{2}\cos(\omega\tau_g)$};
\end{tikzpicture}

  \caption{Two-element interferometer with multiplying correlator. Both
           antennas receive the same signal, delayed by
           $\tau_g = \vec b\cdot\hat s/c$; the averaged product depends
           on the source direction through $\tau_g$. Redrawn after a
           schematic that appears in many publications without a
           traceable original, e.\,g.\ \cite{krinkele2012}.}
  \label{fig:two-element}
\end{figure}
\addtocategory{pic}{krinkele2012}  % listed under figure sources

\section{Complex correlator}

The cosine correlator of Eq.~\eqref{eq:fringe} responds only to the
even part of the brightness distribution. A second correlator with a
$90^\circ$ phase shift in one signal path gives the sine component. Both
together form the \emph{complex visibility} (chapter~\ref{ch:visibility}).
With an SDR receiver both signals are available as complex IQ samples,
so the complex product $x_1 x_2^*$ yields both components directly.

\section{Fringe rate}

The Earth's rotation changes $\tau_g$ continuously. For an east--west
baseline of length $b$ the delay and the fringe frequency are
\begin{equation}
  \tau_g = \frac{b}{c}\cos\delta\,\sin H , \qquad
  f_F = \frac{b}{\lambda}\,\cos\delta\,\cos H\;\omega_E .
  \label{eq:fringerate}
\end{equation}
The fringe rate is largest at transit ($H=0$). For the Sun the solar
rate $\omega_\odot$ must be used instead of the sidereal rate
$\omega_E$; the difference of 0.27\,\% shifts each fringe of a
5-minute period by almost one second.

\section{Bandwidth and delay}

With a finite bandwidth $\Delta\nu$ the fringes of different
frequencies add with different phases. For a rectangular passband the
fringe amplitude is reduced by
\begin{equation}
  \left|\operatorname{sinc}(\Delta\nu\,\tau)\right|, \qquad
  \operatorname{sinc}x = \frac{\sin\pi x}{\pi x},
\end{equation}
where $\tau$ is the residual delay after compensation. Without delay
compensation $\tau = \tau_g$; this limits uncompensated operation to
short baselines (for $\Delta\nu = 2\,\mathrm{MHz}$: loss below 1\,\% up
to $b \approx 10\,\mathrm{m}$). Long baselines need delay tracking.

\section{Sensitivity of an interferometer}

For two antennas observing a point source the correlation coefficient
is
\begin{equation}
  \rho = \frac{\sqrt{T_{A,1}\,T_{A,2}}}
              {\sqrt{(T_{\mathrm{sys},1}+T_{A,1})(T_{\mathrm{sys},2}+T_{A,2})}}
  \;\xrightarrow{\text{identical}}\;
  \frac{T_A}{T_{\mathrm{sys}}+T_A}.
\end{equation}
For a resolved source $\rho$ is multiplied by the normalised
visibility $|\mathcal V|/S$. With Nyquist-sampled signals of bandwidth
$\Delta\nu$ a correlator averages $N = 2\Delta\nu\tau$ independent
products; the noise of each correlator output is $1/\sqrt{N}$ in units
of $\rho$, so
\begin{equation}
  \mathrm{SNR} = \rho\,\sqrt{2\,\Delta\nu\,\tau} .
  \label{eq:snr-corr}
\end{equation}

\section{Adding and multiplying interferometer}

An \emph{adding} interferometer sums the two signals and measures the
total power with a detector. The output is the total-power level plus
the fringe term; the peak-to-peak fringe amplitude relative to the total
power is $2\rho\,|\mathcal V|/S$. A total-power detector is therefore
usable only if the receiver gain is stable to a fraction of this value.
A \emph{multiplying} interferometer outputs only the correlated part and
is insensitive to gain drifts of the uncorrelated noise. Interferosim
models both in \texttt{observation.Correlator}.
